% Options for packages loaded elsewhere
\PassOptionsToPackage{unicode}{hyperref}
\PassOptionsToPackage{hyphens}{url}
\PassOptionsToPackage{dvipsnames,svgnames,x11names}{xcolor}
\documentclass[
  10pt,
  fontset=fandol]{ctexart}
\usepackage{xcolor}
\usepackage[margin=16mm]{geometry}
\usepackage{amsmath,amssymb}
\setcounter{secnumdepth}{-\maxdimen} % remove section numbering
\usepackage{iftex}
\ifPDFTeX
  \usepackage[T1]{fontenc}
  \usepackage[utf8]{inputenc}
  \usepackage{textcomp} % provide euro and other symbols
\else % if luatex or xetex
  \usepackage{unicode-math} % this also loads fontspec
  \defaultfontfeatures{Scale=MatchLowercase}
  \defaultfontfeatures[\rmfamily]{Ligatures=TeX,Scale=1}
\fi
\usepackage{lmodern}
\ifPDFTeX\else
  % xetex/luatex font selection
\fi
% Use upquote if available, for straight quotes in verbatim environments
\IfFileExists{upquote.sty}{\usepackage{upquote}}{}
\IfFileExists{microtype.sty}{% use microtype if available
  \usepackage[]{microtype}
  \UseMicrotypeSet[protrusion]{basicmath} % disable protrusion for tt fonts
}{}
\makeatletter
\@ifundefined{KOMAClassName}{% if non-KOMA class
  \IfFileExists{parskip.sty}{%
    \usepackage{parskip}
  }{% else
    \setlength{\parindent}{0pt}
    \setlength{\parskip}{6pt plus 2pt minus 1pt}}
}{% if KOMA class
  \KOMAoptions{parskip=half}}
\makeatother
\setlength{\emergencystretch}{3em} % prevent overfull lines
\providecommand{\tightlist}{%
  \setlength{\itemsep}{0pt}\setlength{\parskip}{0pt}}
\usepackage{amsmath,amssymb,mathtools,needspace}
\xeCJKDeclareCharClass{CJK}{"2032,"0400->"04FF,"2160->"216B,"2460->"2473,"25A0,"25C7,"25CB,"25CF}
\IfFontExistsTF{Arial Unicode MS}{\xeCJKsetup{AutoFallBack=true}\setCJKfallbackfamilyfont{\CJKrmdefault}{Arial Unicode MS}}{}
\setcounter{secnumdepth}{0}
\setlength{\emergencystretch}{2em}
\allowdisplaybreaks
\usepackage{bookmark}
\IfFileExists{xurl.sty}{\usepackage{xurl}}{} % add URL line breaks if available
\urlstyle{same}
\hypersetup{
  pdftitle={``缀术''再剖析},
  colorlinks=true,
  linkcolor={Maroon},
  filecolor={Maroon},
  citecolor={Blue},
  urlcolor={Blue},
  pdfcreator={LaTeX via pandoc}}

\title{``缀术''再剖析}
\author{}
\date{}

\begin{document}
\maketitle

\subsubsection{12.3.5　``缀术''再剖析}\label{ux7f00ux672fux518dux5256ux6790}

本章一开头指出，南北朝数学家祖冲之给出了高精度的圆周率，这项成就领先世界一千多年。

祖冲之的圆周率是怎样得出来的？如果直接用内接正多边形来逼近，要一直算到正 \(24\,576=6\times2^{12}\) 边形。耗费如此巨大的计算量，在筹算的古代是难以想象的。

中华民族是个智慧的民族，是个善于创新的民族。\textbf{刘徽的加速技术是中华文明前瞻性思维的一个明证。}

据隋唐古书记载，祖冲之的算法设计技术称为``缀术''。史书盛赞缀术``指要精密，

\begin{quote}
校注（agent 补充）：表 12.2 的 \(S_n\) 数值从 3.000 000 000 开始，按原表保留；此前半径 \(r=10\) 寸的表 12.1 数值从 300 开始。本页未明说此处改用单位圆或作了归一化，转录未擅自改变尺度。
\end{quote}

\href{../../source-images/pdf-302.jpeg}{原页图像}

算氏之最''。但关于缀术的具体内容已无从考证，仅仅靠``缀术''这个名词能够窥探出其中的奥秘吗？

查看汉语词典，\textbf{``缀''字有两个含义：一是``缀合''，即组合，因此``缀术''就是组合技术；另一个含义是``缀补''，即修补和校正，在这个意义上，``缀术''亦可理解为校正技术。}

\textbf{这样，所谓``缀术''其实就是组合技术与校正技术。}本节一开头就指出，刘徽的精加工技术正是这种技术。也许，据此可以断定，\textbf{祖冲之的``缀术''正是继承了刘徽的``衣钵''------特别是刘徽的割圆术，归纳总结出来的。}

数学史先辈钱宝琮先生早就猜测过祖冲之与刘徽之间的传承关系。我们深信，\textbf{祖冲之将自己的算法设计技术命名为``缀术''，目的也是试图向世人表白：自己的数学成就，只是前人特别是刘徽的研究工作的缀补和修正。因此，``缀术''实质上只是《九章算术》刘徽注的祖冲之注。}

也许这就是历史的真相。

\begin{center}\rule{0.5\linewidth}{0.5pt}\end{center}

\subsection{相关原页校注（agent 补充）}\label{ux76f8ux5173ux539fux9875ux6821ux6ce8agent-ux8865ux5145}

以下校注来自本节覆盖的原页，可能也涉及同页相邻小节；原文照录与校正说明分开保留。

\subsubsection{原书 PDF 页 302 的校注}\label{ux539fux4e66-pdf-ux9875-302-ux7684ux6821ux6ce8}

\href{../pages/pdf-302.md}{完整原页转录} · \href{../../source-images/pdf-302.jpeg}{原页图像}

校注记录：CH12-E002, CH12-E003, CH12-E004。

\begin{quote}
校注（agent 补充）：本页圆周率末段印作 \(\ldots793\,23\)，PDF294（书页281）同一阿尔·卡西结果印作 \(\ldots792\,23\)，两处原文不一致，均按原页保留。

校注（agent 补充）：原文称 \(L_n=n\sin(\pi/n)\) 为``边长''，该式实际为直径 1 的圆内接正 \(n\) 边形的周长；单边长应为 \(\sin(\pi/n)\)。这是原书文字与所列公式的不一致，正文未作替换。
\end{quote}

\subsection{本节覆盖原页的校注记录}\label{ux672cux8282ux8986ux76d6ux539fux9875ux7684ux6821ux6ce8ux8bb0ux5f55}

下列记录按原页关联，含同页相邻小节的校注；计算时同时读取上文校注和完整原页。

\begin{itemize}
\tightlist
\item
  \texttt{CH12-E002}：原书 PDF 页 294, 302
\item
  \texttt{CH12-E003}：原书 PDF 页 302
\item
  \texttt{CH12-E004}：原书 PDF 页 302, 304
\end{itemize}

\end{document}
